\scanpage{99}
\mspage{49}
le cas au contraire où \deleted{\unsure{il}} $f$ n'est pas définie en $s$, on trouve une \deleted{\unsure{section|seule}} action de $2^{\text{ième}}$ espèce, normalisant un $L_i$ \deleted{\unsure{cyclique}} \unsure{de|dans} $\pi$.

À vrai dire, on n'a regardé jusqu'à présent que ce qui se passe en des ``places'' $s$ de $K/k$ (corps des fonctions relatif d'une variable) où la \deleted{famille} \inserted{courbe} $\unsure{U_\eta/\eta|U_K/K}$ a ``bonne réduction'' --- c'était implicite, car on partait $=$ d'une $U$ sur $S$ --- et je voudrais examiner maintenant ce qui se passe dans le cas où on abandonne cette hypothèse : examiner quels types d'actions de $\mathcal E_{k(s)}\underset{\inserted{\text{ouvert}}}{\subset}\mathcal E_k$ on obtient par la construction précédente, en les $s\in S$ où on a ``mauvaise réduction''. Un cas type serait celui où $k=\mathbf Q$, $K$ corps des fonctions d'une courbe modulaire pour courbes elliptiques (\deleted{\unsure{considérées comme}} \inserted{pointées}), et $s$ place à l'infini de $K$, correspondant à une courbe elliptique
\margin{98}
